%\addtocounter{secnumdepth}{-1}
%\setcounter{secnumdepth}{-1}
\chapter{Resumo}
\label{chap:resumo}

O presente trabalho apresenta o desenvolvimento de um videojogo inovador de simulação e gestão estratégica de infraestruturas tecnológicas, que combina a complexidade e o rigor técnico do ecossistema de computação com uma experiência lúdica altamente envolvente. Ao assumir o papel de gestor de um \textit{data center}, o jogador é desafiado a garantir a viabilidade económica da sua operação através da otimização de recursos financeiros, alocação de energia e seleção de componentes de \textit{hardware} reais.


Desenvolvido no motor Unity e aliado a metodologias sólidas de game design, como os \textit{frameworks} \ac{MDA}, \textit{Flow} e o Modelo de Hook, o produto destaca-se por transformar o conhecimento técnico em ferramentas essenciais de progressão. A inclusão de mecânicas dinâmicas de combate contra ameaças cibernéticas e de \textit{minigames} baseados em terminais virtuais (\textit{WordRush}) permite quebrar a monotonia da simulação económica, garantindo um ritmo contínuo de adrenalina e retenção do jogador.

Testado e avaliado com utilizadores reais, o videojogo demonstrou um elevado potencial na transmissão indireta de conceitos avançados de \ac{TI} e cibersegurança, traduzindo-se numa solução de entretenimento educativo pronta a capturar o interesse de entusiastas de jogos de estratégia, simulação e tecnologia.

\section*{Palavras-chave}
Jogo Educacional, Jogos Sérios, Jogo de Matemática, Unity